\section{Классификация квадратичных форм в вещественном пространстве. Критерий знакоопределенности квадратичных форм}

\begin{definition}
Квадратичная форма $f(x, x)$ над $\mathbb{R}$ называется:
\begin{enumerate}
    \item \textbf{Положительно определенной}, если $\forall x \neq 0 \implies f(x, x) > 0$.
    \item \textbf{Отрицательно определенной}, если $\forall x \neq 0 \implies f(x, x) < 0$.
    \item \textbf{Неотрицательно определенной} (положительно полуопределенной), если $\forall x \implies f(x, x) \ge 0$, и $\exists x \neq 0 \implies f(x, x) = 0$.
    \item \textbf{Неположительно определенной} (отрицательно полуопределенной), если $\forall x \implies f(x, x) \le 0$, и $\exists x \neq 0 \implies f(x, x) = 0$.
    \item \textbf{Незнакоопределенной}, если $\exists x_1, x_2$ такие, что $f(x_1, x_1) > 0$ и $f(x_2, x_2) < 0$.
\end{enumerate}
\end{definition}

\begin{theorem}[Классификация через индексы инерции]
Пусть $p$ и $q$ --- индексы инерции квадратичной формы $f$ в $n$-мерном пространстве. Тогда:
\begin{enumerate}
    \item $f$ положительно определена $\iff p = n$ (и $q = 0$).
    \item $f$ отрицательно определена $\iff q = n$ (и $p = 0$).
    \item $f$ неотрицательно определена $\iff q = 0$ и $p < n$.
    \item $f$ неположительно определена $\iff p = 0$ и $q < n$.
    \item $f$ незнакоопределена $\iff p > 0$ и $q > 0$.
\end{enumerate}
\end{theorem}

\begin{proof}
Следует непосредственно из нормального вида формы $f(x, x) = \xi_1^2 + \dots + \xi_p^2 - \xi_{p+1}^2 - \dots - \xi_{p+q}^2$. Например, если $p=n$, то $f(x, x) = \sum_{i=1}^n \xi_i^2$, что строго больше нуля для любого $x \neq 0$. Если $p>0$ и $q>0$, то, взяв вектор с $\xi_1=1$ (остальные 0), получим $f>0$, а с $\xi_{p+1}=1$ (остальные 0) --- $f<0$, что дает незнакоопределенность.
\hfill$\blacksquare$
\end{proof}

\begin{definition}
Пусть $A$ --- симметричная матрица квадратичной формы. \textbf{Главными минорами} $\Delta_k$ матрицы $A$ называются определители её верхних левых квадратных подматриц порядка $k$:
\[
\Delta_1 = a_{11}, \quad 
\Delta_2 = \begin{vmatrix} a_{11} & a_{12} \\ a_{21} & a_{22} \end{vmatrix}, \quad \dots, \quad 
\Delta_n = \det A = \begin{vmatrix} 
a_{11} & a_{12} & \dots & a_{1n} \\ 
a_{21} & a_{22} & \dots & a_{2n} \\ 
\vdots & \vdots & \ddots & \vdots \\ 
a_{n1} & a_{n2} & \dots & a_{nn} 
\end{vmatrix}.
\]
\end{definition}

\begin{theorem}[Критерий Сильвестра]
Квадратичная форма $f(x, x)$ с матрицей $A$ в некотором базисе является:
\begin{enumerate}
    \item \textbf{Положительно определенной} тогда и только тогда, когда все главные миноры строго положительны:
    \[
    \Delta_1 > 0, \quad \Delta_2 > 0, \quad \dots, \quad \Delta_n > 0.
    \]
    \item \textbf{Отрицательно определенной} тогда и только тогда, когда главные миноры знакопеременны, начиная с отрицательного:
    \[
    \Delta_1 < 0, \quad \Delta_2 > 0, \quad \Delta_3 < 0, \quad \dots, \quad (-1)^n \Delta_n > 0.
    \]
\end{enumerate}
\end{theorem}

\begin{proof}
\begin{enumerate}
    \item \textbf{Необходимость (для положительной определенности):} 
    Предположим, что $f(x, x) > 0$ для всех $x \neq 0$, но существует главный минор $\Delta_k = 0$. Рассмотрим однородную систему линейных уравнений, составленную из первых $k$ уравнений системы $Ax = 0$:
    \[
    \begin{cases}
    a_{11}x_1 + \dots + a_{1k}x_k = 0 \\
    \vdots \\
    a_{k1}x_1 + \dots + a_{kk}x_k = 0
    \end{cases}
    \]
    Определитель этой системы равен $\Delta_k = 0$, следовательно, она имеет нетривиальное решение $x_0 = (x_1^0, \dots, x_k^0, 0, \dots, 0)^T \neq 0$. 
    Для этого вектора значение формы равно:
    \[
    f(x_0, x_0) = \sum_{i,j=1}^k a_{ij} x_i^0 x_j^0 = x_0^T A x_0 = 0,
    \]
    так как первые $k$ уравнений выполнены, а остальные координаты $x_0$ равны нулю. Это противоречит положительной определенности формы ($f(x_0, x_0)$ должно быть $>0$). 
    Следовательно, все $\Delta_k \neq 0$. Поскольку форма положительно определена, по методу Якоби (применимому, так как все $\Delta_k \neq 0$) она приводится к виду $f(x, x) = \sum_{i=1}^n \frac{\Delta_i}{\Delta_{i-1}} \eta_i^2$ (где $\Delta_0 = 1$). Чтобы форма была положительной, все коэффициенты $\frac{\Delta_i}{\Delta_{i-1}}$ должны быть положительны. Отсюда индукцией получаем $\Delta_1 > 0, \Delta_2 > 0, \dots, \Delta_n > 0$.

    \item \textbf{Достаточность:} 
    Если все $\Delta_k > 0$, то метод Якоби применим, и форма приводится к каноническому виду $f(x, x) = \sum_{i=1}^n \frac{\Delta_i}{\Delta_{i-1}} \eta_i^2$. Поскольку все $\Delta_i > 0$, все коэффициенты при квадратах строго положительны. Следовательно, для любого $x \neq 0$ (а значит, и $\eta \neq 0$) сумма квадратов с положительными коэффициентами строго больше нуля. Форма положительно определена.

    \item \textbf{Отрицательная определенность:} 
    Форма $f(x, x)$ отрицательно определена тогда и только тогда, когда форма $-f(x, x)$ положительно определена. Матрица формы $-f$ равна $-A$. Её главные миноры равны $(-1)^k \Delta_k$. Применяя к $-f$ критерий положительной определенности, получаем условие $(-1)^k \Delta_k > 0$ для всех $k$, что эквивалентно знакопеременности миноров $\Delta_k$, начиная с отрицательного.
\end{enumerate}
\hfill$\blacksquare$
\end{proof}
